\documentclass[11pt,letterpaper]{article}

\usepackage[top=1in, bottom=1in, left=1in, right=1in, headheight=14pt]{geometry}
\usepackage{fontspec}
\setmainfont{DejaVu Serif}
\setsansfont{DejaVu Sans}
\setmonofont{DejaVu Sans Mono}

\usepackage{xcolor}
\usepackage{titlesec}
\usepackage{fancyhdr}
\usepackage{enumitem}
\usepackage{hyperref}
\usepackage{booktabs}
\usepackage{tabularx}
\usepackage{longtable}
\usepackage{colortbl}
\usepackage{multirow}
\usepackage{array}
\usepackage{parskip}
\usepackage{tcolorbox}
\usepackage{graphicx}
\usepackage{lastpage}

% Color Scheme: Ecology/Nature
\definecolor{primary}{HTML}{1B5E20}
\definecolor{secondary}{HTML}{1565C0}
\definecolor{tertiary}{HTML}{4E342E}
\definecolor{accent}{HTML}{2E7D32}
\definecolor{lightprimary}{HTML}{E8F5E9}
\definecolor{lightsecondary}{HTML}{E3F2FD}
\definecolor{lighttertiary}{HTML}{EFEBE9}
\definecolor{rowgray}{HTML}{F5F5F5}
\definecolor{bordergray}{HTML}{BDBDBD}
\definecolor{subtlegray}{HTML}{757575}
\definecolor{darktext}{HTML}{212121}

% Section Formatting
\titleformat{\section}
  {\Large\bfseries\sffamily\color{primary}}
  {\thesection}{1em}{}
  [\vspace{-0.5em}\textcolor{bordergray}{\rule{\textwidth}{0.4pt}}]

\titleformat{\subsection}
  {\large\bfseries\sffamily\color{secondary}}
  {\thesubsection}{1em}{}

\titleformat{\subsubsection}
  {\normalsize\bfseries\sffamily\color{tertiary}}
  {\thesubsubsection}{1em}{}

\titlespacing*{\section}{0pt}{1.5em}{0.8em}
\titlespacing*{\subsection}{0pt}{1.2em}{0.5em}
\titlespacing*{\subsubsection}{0pt}{0.8em}{0.3em}

% Custom Environments
\newcommand{\stageheader}[2]{%
  \noindent\colorbox{#2}{\makebox[\textwidth][l]{%
    \textcolor{white}{\bfseries\sffamily\hspace{6pt}#1\hspace{6pt}}}}%
  \vspace{0.5em}\par%
}

\newtcolorbox{asciibox}{
  colback=rowgray, colframe=bordergray,
  arc=1pt, boxrule=0.5pt,
  left=6pt, right=6pt, top=4pt, bottom=4pt,
  fontupper=\small\ttfamily,
  breakable
}

\newtcolorbox{quotebox}{
  colback=lightprimary, colframe=primary,
  arc=2pt, boxrule=0.5pt,
  left=12pt, right=12pt, top=8pt, bottom=8pt,
  breakable
}

\newcommand{\metafield}[2]{\textbf{\sffamily #1} #2\par}
\newcolumntype{L}[1]{>{\raggedright\arraybackslash}p{#1}}
\newcolumntype{C}[1]{>{\centering\arraybackslash}p{#1}}
\newcommand{\tableheader}[1]{\textbf{\sffamily\textcolor{white}{#1}}}

% Headers and Footers
\pagestyle{fancy}
\fancyhf{}
\fancyhead[L]{\small\sffamily\textcolor{subtlegray}{Pacific Northwest Gardening}}
\fancyhead[R]{\small\sffamily\textcolor{subtlegray}{GSD Mission Package}}
\fancyfoot[C]{\small\sffamily\textcolor{subtlegray}{Page \thepage\ of \pageref{LastPage}}}
\renewcommand{\headrulewidth}{0.4pt}
\renewcommand{\footrulewidth}{0pt}

% Hyperlinks
\hypersetup{
  colorlinks=true,
  linkcolor=secondary,
  urlcolor=secondary,
  pdfauthor={GSD Ecosystem},
  pdftitle={Pacific Northwest Gardening -- Mission Package},
  pdfsubject={Vision to Mission Pipeline},
}

\begin{document}

% TITLE PAGE
\thispagestyle{empty}
\begin{center}
\vspace*{3cm}

{\small\sffamily\textcolor{primary}{\textbf{GSD RESEARCH MISSION PACKAGE}}}

\vspace{0.3cm}
\textcolor{bordergray}{\rule{8cm}{0.4pt}}

\vspace{1.5cm}
{\fontsize{28}{34}\selectfont\bfseries\sffamily\textcolor{primary}{Gardening in the\\[0.3em]Pacific Northwest}}

\vspace{0.8cm}
{\Large\sffamily\textcolor{secondary}{Western Washington, Oregon \& Coastal British Columbia}}

\vspace{0.5cm}
{\large\sffamily\itshape\textcolor{tertiary}{A Deep Research Study}}

\vspace{3cm}
{\sffamily\textcolor{accent}{Vision $\rightarrow$ Research $\rightarrow$ Mission}}\\[0.3em]
{\small\sffamily\textcolor{accent}{Full Pipeline Execution}}

\vspace{3cm}
{\small\sffamily\textcolor{subtlegray}{Date: March 7, 2026  |  Status: Ready for GSD Orchestrator}}\\[0.3em]
{\small\sffamily\textcolor{subtlegray}{Activation Profile: Squadron  |  Estimated: 12 context windows across 4 sessions}}
\end{center}
\newpage

% TABLE OF CONTENTS
\tableofcontents
\newpage

% STAGE 1 --- VISION DOCUMENT
\stageheader{STAGE 1 --- VISION DOCUMENT}{primary}

\section{Pacific Northwest Gardening --- Vision Guide}

\metafield{Date:}{March 7, 2026}
\metafield{Status:}{Initial Vision / Pre-Research}
\metafield{Depends on:}{None (root document)}
\metafield{Context:}{GSD Ecosystem --- Ecology/Education Domain}

\subsection{Vision}

The Pacific Northwest occupies a paradox in the gardening world. It is simultaneously one of the most productive growing regions on the planet and one of the most misunderstood. Gardeners arriving from other parts of North America bring assumptions forged in continental climates---expectations about when to plant, what will thrive, and how soil behaves---that quietly sabotage their first seasons. The region's maritime influence, its volcanic soils, its extraordinary native plant heritage, and its rapidly shifting hardiness zones under climate change all demand a fundamentally different approach.

This research mission aims to produce a comprehensive, authoritative reference for gardening in the Pacific Northwest---one grounded in university extension research, government agency data, and the accumulated wisdom of the region's professional horticulturists. The result will be a living document that serves both the beginning gardener planting their first raised bed in a Seattle backyard and the experienced grower adapting to hardiness zone shifts that now place the Puget Sound lowlands in USDA Zone 9a.

The deeper ambition is to model how a garden functions as an ecosystem rather than merely a production facility. Native plant integration, pollinator support, soil biology, and climate adaptation are not optional add-ons---they are the architecture of a resilient garden. In the GSD philosophy, this mirrors the Amiga Principle: achieving remarkable outcomes through ecological leverage rather than brute-force inputs of water, fertilizer, and pesticide.

The spaces between---between seasons, between native and cultivated, between the garden and the wild landscape that surrounds it---are where the real intelligence of Pacific Northwest gardening lives.

\subsection{Problem Statement}

\begin{enumerate}[leftmargin=2em]
  \item \textbf{Climate Mismatch:} Most general gardening guides are written for continental climates. PNW gardeners face cool wet springs, dry summers, mild wet winters, and acidic soils---a Mediterranean-maritime pattern that requires fundamentally different timing, plant selection, and soil management strategies.

  \item \textbf{Shifting Hardiness Zones:} The 2023 USDA Hardiness Zone Map moved much of the Puget Sound lowlands from Zone 8b to 9a, reflecting a five-degree increase in average minimum winter temperatures. Gardeners need guidance on what this shift means in practice---which plants can now overwinter, which traditional assumptions no longer hold, and how to plan for continued change.

  \item \textbf{Native Plant Disconnection:} Home landscapes in the PNW frequently rely on non-native ornamentals that provide little ecological value. Meanwhile, the region's extraordinary native flora---adapted over millennia to local soils and climate---remains underutilized in residential gardens, leaving gaps in pollinator habitat and biodiversity support.

  \item \textbf{Soil Complexity and Acidity:} PNW soils are a patchwork of glacial deposits, volcanic ash, heavy clay, and sandy loam. High rainfall produces acidic conditions (often below pH 6.0) that lock out nutrients. Without understanding local soil composition, gardeners waste effort and resources on amendments that may do more harm than good.

  \item \textbf{Pest and Disease Pressure Under Moisture:} The region's cool, humid conditions create ideal habitat for slugs, fungal diseases (powdery mildew, apple scab, anthracnose), and specific insect pests (imported cabbageworm, black vine weevil, crane fly larvae) that demand integrated, ecologically-aware management rather than broad-spectrum chemical intervention.
\end{enumerate}

\subsection{Core Concept}

\begin{quotebox}
\textbf{One-line model:} The Pacific Northwest garden is a managed ecosystem operating within a maritime-Mediterranean climate gradient, where success is measured by ecological integration rather than resistance to natural conditions.
\end{quotebox}

The core concept reframes gardening from ``conquering'' the PNW climate to collaborating with it. This means embracing the region's natural rhythms: cool-season crops in spring and fall, drought-adapted natives in summer, soil biology that thrives under no-till management, and year-round food production that exploits the mild winter rather than shutting down during it.

\subsubsection{Study Region}

\begin{asciibox}
\begin{verbatim}
                    PACIFIC NORTHWEST GARDENING REGION
    +------------------------------------------------------+
    |  Vancouver, BC --- Fraser Valley                     |
    |       |                                              |
    |  Puget Sound ---- San Juan Islands                   |
    |  (Zone 9a)        |                                  |
    |       |      Olympic Peninsula                       |
    |  Portland ---- Willamette Valley                     |
    |  (Zone 8b-9a)     |                                  |
    |       |      Oregon Coast                            |
    |  Roseburg ---- Umpqua Valley                         |
    |                                                      |
    |  === CASCADE RANGE (climatic divide) ===             |
    |                                                      |
    |  East: Semi-arid, continental (out of scope v1.0)    |
    +------------------------------------------------------+
    Focus: West of the Cascades, sea level to ~1500 ft
    Zones: USDA 7b through 9a | Sunset Zones 4-6
\end{verbatim}
\end{asciibox}

\subsubsection{Module Architecture}

\begin{asciibox}
\begin{verbatim}
    +------------------+     +-------------------+
    |  M1: CLIMATE &   |---->|  M2: SOIL SCIENCE |
    |  MICROCLIMATES   |     |  & MANAGEMENT     |
    +--------+---------+     +--------+----------+
             |                        |
             v                        v
    +------------------+     +-------------------+
    |  M3: NATIVE      |<--->|  M4: FOOD         |
    |  PLANT           |     |  PRODUCTION &     |
    |  INTEGRATION     |     |  SEASONAL CALENDAR|
    +--------+---------+     +--------+----------+
             |                        |
             +----------+-------------+
                        v
             +-------------------+
             |  M5: PEST,       |
             |  DISEASE &       |
             |  CLIMATE         |
             |  ADAPTATION      |
             +-------------------+
\end{verbatim}
\end{asciibox}

\subsection{Research Modules}

\subsubsection{M1: Climate and Microclimates}

USDA hardiness zones (7b--9a), Sunset climate zones, the maritime-Mediterranean pattern, rain shadow effects, the Cascade divide, El Nino/La Nina oscillation impacts, frost date ranges, heat unit accumulation, and the critical concept of microclimates---south-facing walls as heat batteries, cold-air drainage in valleys, coastal buffering. Includes the 2023 zone map shift and its practical implications.

\subsubsection{M2: Soil Science and Management}

Glacial-volcanic soil origins, clay/silt/sand/loam distribution across the region, soil pH and acidity management (lime application thresholds), organic matter building through compost and mulch, no-till and sheet mulching techniques, cover cropping for winter, soil testing protocols (King Conservation District, WSU Extension, OSU Extension), raised bed strategies for drainage, and the critical role of soil biology (mycorrhizal networks, beneficial microorganisms).

\subsubsection{M3: Native Plant Integration}

The PNW native plant palette: trees (Douglas-fir, western redcedar, Oregon white oak, vine maple), shrubs (Oregon grape, salal, red-flowering currant, oceanspray, elderberry, evergreen huckleberry), groundcovers (kinnikinnick, coastal strawberry, sword fern), and wildflowers (camas, trillium, bleeding heart, yarrow). Ecological functions: pollinator support, wildlife habitat, erosion control, water-wise landscaping. Plant community associations and companion planting with natives. Backyard Habitat Certification programs. Local genotype sourcing.

\subsubsection{M4: Food Production and Seasonal Calendar}

Month-by-month planting calendar for the maritime PNW (Vancouver BC to Roseburg OR). Cool-season crop emphasis (brassicas, greens, root vegetables, alliums). Warm-season strategies (short-season tomato varieties, heat-accumulation techniques). The ``second spring'' of August---fall and winter crop establishment. Year-round harvest potential. Succession planting. Raised bed construction and management. Greenhouse and cold frame season extension. Key reference texts from Tilth Alliance and regional growers.

\subsubsection{M5: Pest, Disease, and Climate Adaptation}

Slug and snail management (iron phosphate baits, habitat modification, beneficial predators). Imported cabbageworm and row cover strategies. Fungal disease prevention (powdery mildew, apple scab, anthracnose)---cultivar selection, airflow, timing. Emerging threats: emerald ash borer (first detected Forest Grove, OR, June 2022), chafer beetle larvae. Integrated pest management philosophy. Climate adaptation: drought-resistant varieties, water-wise irrigation (drip, soaker hose), shifting planting windows, heat dome preparedness (shade cloth, evaporative cooling), no-till for soil moisture retention.

\subsection{Chipset Configuration}

\begin{asciibox}
\begin{verbatim}
chipset:
  name: pnw-gardening-research
  version: 1.0
  activation: squadron
  domain: ecology-education

  modules:
    - id: M1-climate
      name: "Climate & Microclimates"
      model: sonnet
      priority: foundation

    - id: M2-soil
      name: "Soil Science & Management"
      model: sonnet
      priority: core

    - id: M3-native
      name: "Native Plant Integration"
      model: opus
      priority: core

    - id: M4-food
      name: "Food Production & Seasonal Calendar"
      model: sonnet
      priority: core

    - id: M5-pest-adapt
      name: "Pest, Disease & Climate Adaptation"
      model: opus
      priority: core

  cross-cutting:
    - safety: source-quality-enforcement
    - integration: ecosystem-connectivity
    - verification: citation-traceability
\end{verbatim}
\end{asciibox}

\subsection{Scope Boundaries}

\textbf{In Scope (v1.0):}
\begin{itemize}[leftmargin=2em]
  \item Maritime PNW west of the Cascades, from coastal BC to southern Oregon
  \item USDA zones 7b through 9a; Sunset zones 4--6
  \item Residential-scale gardening (backyard, community garden, small urban farm)
  \item Vegetable production, native plant landscaping, ornamental gardening
  \item Climate adaptation strategies for home gardeners
  \item Integrated pest management using organic and ecological methods
\end{itemize}

\textbf{Out of Scope:}
\begin{itemize}[leftmargin=2em]
  \item East of Cascades (semi-arid, continental climate---separate mission)
  \item Commercial-scale agriculture and row crop farming
  \item Indoor cannabis or controlled-environment agriculture
  \item Detailed landscape architecture or hardscape design
  \item Forestry and timber management
\end{itemize}

\subsection{Success Criteria}

\begin{enumerate}[leftmargin=2em]
  \item Climate module covers all three PNW zone systems (USDA, Sunset, AHS Heat) with practical interpretation for gardeners.
  \item Soil module addresses all four major PNW soil types (clay, silt, sand, loam) with amendment protocols sourced from university extension.
  \item Native plant inventory includes at minimum 25 species across trees, shrubs, groundcovers, and wildflowers, each with habitat requirements and ecological function.
  \item Food production calendar provides month-by-month guidance for at least 20 common vegetables with region-specific timing.
  \item Pest and disease module covers the five most common PNW garden pests with organic/IPM management strategies.
  \item All climate projections cite USDA Climate Hubs, NOAA, or peer-reviewed research.
  \item Every plant recommendation includes USDA hardiness zone range and sun/moisture requirements.
  \item Climate adaptation section addresses drought, heat dome events, and shifting frost dates with practical garden-level responses.
  \item Cross-module integration: native plants are connected to pest management (beneficial insect habitat) and soil health (mycorrhizal networks).
  \item All numerical claims (temperatures, pH ranges, zone designations, pest counts) attributed to specific sources.
  \item Document is self-contained---a reader with no prior PNW gardening knowledge can understand and act on all recommendations.
  \item Safety-critical: no policy advocacy; evidence-based presentation throughout.
\end{enumerate}

\subsection{Relationship to Other Documents}

\begin{center}
\rowcolors{2}{rowgray}{white}
\begin{tabularx}{\textwidth}{L{4cm}XC{2.5cm}}
\toprule
\rowcolor{primary}
\tableheader{Document} & \tableheader{Relationship} & \tableheader{Type} \\
\midrule
Heritage Skills Pack & Cultural plant knowledge, Foxfire traditions & Upstream \\
Cooking with Claude & Recipe integration, garden-to-table pipeline & Downstream \\
Dependency Health Monitor & Package freshness model applied to seed/cultivar data & Pattern \\
GSD Skill Creator & Execution engine for this mission & Infrastructure \\
The Space Between & Mathematical foundations; ecosystem as complex system & Philosophical \\
\bottomrule
\end{tabularx}
\end{center}

\subsection{The Through-Line}

A Pacific Northwest garden, tended well, is a living expression of the Amiga Principle. It achieves extraordinary productivity not through overwhelming inputs but through architectural leverage---the right plant in the right place at the right time, working with the region's natural rhythms rather than against them. The spaces between seasons, between native and cultivated, between the garden bed and the forest edge---these are where resilience lives.

In this way, a garden mirrors the GSD ecosystem itself: specialized agents (plants, soil organisms, pollinators) operating within a structured environment (climate, soil chemistry, seasonal cycles), each contributing to an emergent whole that exceeds the sum of its parts. The gardener, like the orchestrator, does not command nature but creates the conditions under which nature's intelligence can express itself.

This is technology in its oldest and most humane form: the deliberate cultivation of living systems for mutual benefit. It is the work that gives people their lives back---hands in soil, family around a table, the giddy satisfaction of a first ripe tomato still warm from the sun.

\newpage

% STAGE 2 --- RESEARCH REFERENCE
\stageheader{STAGE 2 --- RESEARCH REFERENCE}{secondary}

\section{Pacific Northwest Gardening --- Research Reference}

\subsection{How to Use This Document}

This reference compiles findings from government agencies, university extension services, and professional horticultural organizations. Each module section provides the factual foundation that mission agents will use to produce final deliverables. Key source organizations include:

\begin{itemize}[leftmargin=2em]
  \item USDA Plant Hardiness Zone Map (2023 revision, with Oregon State University PRISM Climate Group)
  \item Oregon State University Extension Service (EC-1577, monthly garden calendars, native plant guides)
  \item Washington State University Extension (EM051E, Hortsense/Pestsense databases, PNW Vegetable Extension Group)
  \item USDA Northwest Climate Hub (drought impacts, economic analysis, climate projections)
  \item U.S. EPA Climate Impacts reports (Northwest regional assessment)
  \item Tilth Alliance (Maritime Northwest Garden Guide, organic gardening standards)
  \item King County Department of Natural Resources (native plant resources)
  \item Pacific Northwest Pest Management Handbooks (OSU / WSU / U of Idaho joint publication)
  \item Washington Native Plant Society
  \item Portland Bureau of Parks (tree disease and pest identification)
\end{itemize}

\subsection{M1 Reference: Climate and Microclimates}

\subsubsection{USDA Hardiness Zones in the PNW}

The Pacific Northwest spans USDA Hardiness Zones 4 (mountain elevations) through 9a (Puget Sound and Willamette Valley lowlands). The critical gardening zones west of the Cascades range from 7b to 9a. The 2023 USDA map update, developed with Oregon State University's PRISM Climate Group, incorporated data from 13,412 weather stations (compared to 7,983 in the 2012 map) and shifted many lowland areas warmer by one half-zone.

The Shoreline/North Seattle area moved from Zone 8b to 9a, representing a five-degree (Fahrenheit) increase in the 30-year average annual extreme minimum temperature. The USDA clarifies that zone designations are based on 30-year averaged data and reflect both improved measurement technology and actual warming trends.

\subsubsection{Maritime-Mediterranean Climate Pattern}

Western PNW climate is characterized by cool, wet winters (October through April) and warm, dry summers (June through September). Spring arrives gradually with persistent cool, overcast conditions that delay soil warming. Autumn is long and gentle, with frost dates varying from early October at higher elevations to late November near the coast.

Key climate metrics for western Washington: average annual temperature ranges from 51 degrees F at the coast to 40 degrees F in the northeast mountains (USDA/Gilmour data). Winter weather is driven by Pacific low-pressure cyclone systems bringing moist, warm air; summer is governed by a high-pressure anticyclone system bringing cool air and dry conditions.

\subsubsection{Microclimates and Heat Accumulation}

South-facing walls, fences, and slopes create warmer microclimates that act as heat batteries, absorbing warmth during the day and radiating it to adjacent plants at night. Low-lying areas collect cold air and experience later frost dates. Coastal areas are buffered from temperature extremes but receive more wind and salt spray.

The concept of ``climate analog'' gardening---sourcing plants from regions worldwide with similar elevation, soil, and sun exposure---is gaining traction. The University of Washington Arboretum has begun trialing crepe myrtles from the Southeast and manzanitas from northwestern California and southern Oregon highlands, targeting species from higher elevations that tolerate both cold and dry summer heat.

\subsection{M2 Reference: Soil Science and Management}

\subsubsection{Soil Origins and Types}

PNW soils west of the Cascades are shaped by volcanic activity, glacial deposits, and millennia of forest cover. The resulting soil landscape is highly variable, often changing character within a single lot. Primary soil types (WSU Extension, NRCS):

\textbf{Heavy clay:} Dense, sticky when wet, rock-hard when dry. Drains slowly, compacts easily, but typically high in nutrients. Requires organic matter and drainage planning. Common in the Puget Sound lowlands.

\textbf{Sandy loam:} Loose, well-draining, easy to dig. Excellent for root vegetables and Mediterranean plants but poor at retaining water and nutrients. Found in glacial outwash areas.

\textbf{Silty soil:} Fine-textured, smooth, holds water well but compacts readily. Mixing in coarse materials improves structure. Common in river valleys and floodplains.

\textbf{Mixed glacial deposits:} Many PNW yards contain variable soils---potentially clay, sand, gravel, and silt within the same garden, requiring zone-by-zone assessment.

\subsubsection{Soil pH and Acidity Management}

High rainfall leaches calcium and magnesium from PNW soils, producing acidic conditions. Many garden soils test below pH 6.0 (St. Clare Heirloom Seeds, OSU Extension data). Soil testing is available through King Conservation District (free for eligible landowners), WSU Extension, and OSU Extension.

When pH falls below 6.0, adding garden lime (calcium carbonate) in spring helps vegetables access soil nutrients more effectively. Overapplication of compost can be problematic---soil as dark as coffee grounds indicates potential nutrient overload, in which case a coarse-textured surface mulch (2--4 inches of wood chips or bark) is more appropriate (Pacific Horticulture Society).

\subsubsection{No-Till and Soil Biology}

Shifting to no-till approaches---layering organic matter on the soil surface rather than tilling---preserves beneficial soil organisms and mimics natural forest-floor processes (Pacific Horticulture Society, Tilth Alliance). Cover crops planted in fall protect soil from winter erosion and add organic matter. Sheet mulching converts lawn to garden bed without tilling. Horticultural coir is recommended over peat moss for soil amendment due to root-stimulating hormones and modest nutrient levels (Pacific Horticulture Society).

\subsection{M3 Reference: Native Plant Integration}

\subsubsection{Ecological Function of Native Plants}

PNW native plants have evolved in the region's wet-winter, dry-summer climate for thousands of years (OSU Extension EC-1577, Washington Native Plant Society). Key benefits: require less water once established, resist native pests and diseases, need less fertilizer and no pesticides, and support pollinator populations that non-native landscapes cannot. As wild spaces decrease, adding native plants to urban landscapes becomes increasingly important for regional biodiversity (Portland Nursery, King County DNRP).

The most important native host plants for moths and butterflies in the PNW---by abundance of species supported---include Oregon white oak, Douglas-fir, willows, and native cherry species (Real Gardens Grow Natives).

\subsubsection{Key Native Plant Palette}

\textbf{Canopy trees:} Douglas-fir (\textit{Pseudotsuga menziesii}), western redcedar (\textit{Thuja plicata}), bigleaf maple (\textit{Acer macrophyllum}), Oregon white oak (\textit{Quercus garryana}---requires dry conditions, subject to root rot if irrigated; OSU Extension).

\textbf{Understory trees and shrubs:} Vine maple (\textit{Acer circinatum}), tall Oregon grape (\textit{Mahonia aquifolium}---Zone 5, sun or shade), red-flowering currant (\textit{Ribes sanguineum}---year-round appeal, wildlife food), Pacific wax myrtle (\textit{Morella californica}---Zone 7, coastal/poor soil tolerant), oceanspray (\textit{Holodiscus discolor}), native red elderberry (\textit{Sambucus racemosa}), evergreen huckleberry (\textit{Vaccinium ovatum}---Zone 7), red huckleberry (\textit{Vaccinium parvifolium}).

\textbf{Groundcovers and ferns:} Sword fern (\textit{Polystichum munitum}---4--6 ft, sun or shade, erosion control), kinnikinnick (\textit{Arctostaphylos uva-ursi}---drought-tolerant, sunny well-drained areas), coastal strawberry (\textit{Fragaria chiloensis}---drought-tolerant groundcover; OSU Extension), maidenhair fern (\textit{Adiantum pedatum}---shade, moist acidic soil, deer resistant), salal (\textit{Gaultheria shallon}---robust, edible berries, shady acidic soil).

\textbf{Wildflowers and perennials:} Common camas (\textit{Camassia quamash}---requires summer dry periods, drought-tolerant, full sun), western trillium (\textit{Trillium ovatum}---spring, partial shade), Pacific bleeding heart (\textit{Dicentra formosa}---hummingbird-attracting, moist shade, dormant in summer), yarrow (\textit{Achillea millefolium}---drought-tolerant, full sun, pollinator magnet), showy milkweed (\textit{Asclepias speciosa}---pollinator attractant).

\subsubsection{Plant Community Design}

Native plants perform best in natural community associations. Red-flowering currant associates with Douglas-fir, bigleaf maple, madrone, vine maple, oceanspray, mock orange, serviceberry, salal, sword fern, and kinnikinnick (Real Gardens Grow Natives). Sourcing plants from local genotypes preserves genetic diversity critical for climate adaptation.

Portland Audubon and Columbia Land Trust operate a Backyard Habitat Certification program guiding homeowners through native plant integration for wildlife support (Portland Nursery).

\subsection{M4 Reference: Food Production and Seasonal Calendar}

\subsubsection{Year-Round Harvest Potential}

The maritime PNW's mild winters enable year-round vegetable production, a capability unmatched in most of North America (Tilth Alliance). Success depends on three strategic shifts: embracing cool-season crops as the primary production window, using the August ``second spring'' for fall/winter planting, and managing soil drainage rather than soil moisture (St. Clare Heirloom Seeds).

\subsubsection{Seasonal Framework}

\textbf{January--February:} Soil testing, seed ordering, pruning dormant fruit trees. Start cool-season seeds indoors. In mild years, direct-sow peas and fava beans outdoors in late February.

\textbf{March--April:} Direct-sow carrots, beets, radishes, spinach, lettuce, and chard as soil reaches 40 degrees F. Transplant brassica starts. Prepare raised beds. Apply lime if needed. Plant potatoes in late March (OSU Extension monthly calendars).

\textbf{May--June:} Transplant warm-season crops after soil reaches 60 degrees F and nights stay above 50 degrees F---typically mid-May to early June. Common mistake: planting tomatoes during a warm early-May week, only to have cool weather stunt them (St. Clare Heirloom Seeds). Use short-season cultivars. Succession-plant beans, lettuce, root crops.

\textbf{July--August:} Peak warm-season harvest. Begin the ``second spring''---sow fall kale, carrots, lettuce, spinach, beets in August. These grow through autumn and can be harvested into December (St. Clare Heirloom Seeds, Deep Harvest Farm). Irrigate consistently during the dry season.

\textbf{September--October:} Continue fall harvest. Plant garlic in October. Sow overwintering cover crops. Begin mulching beds.

\textbf{November--December:} Harvest frost-sweetened kale, leeks, Brussels sprouts, winter greens. Monitor drainage and slug activity. Plan next season.

\subsubsection{Top Performing Vegetables}

Well-suited to PNW conditions: beans, beets, broccoli, carrots, cauliflower, chard, garlic, kale, leeks, lettuce, onions, peas, potatoes, radishes, spinach, and squash (ParentMap, WSU Extension). Crops that struggle without season extension: eggplant, peppers, full-season corn and watermelon---short-season cultivars are essential.

\subsubsection{Season Extension}

Raised beds solve the two biggest PNW problems: poor drainage and cold soil temperatures (St. Clare Heirloom Seeds). Cold frames and low tunnels add 4--6 weeks to both ends of the season. Greenhouses extend growing to 8--9 months inland and potentially year-round near the coast (Planta Greenhouses, WSU Extension).

\subsection{M5 Reference: Pest, Disease, and Climate Adaptation}

\subsubsection{Slugs and Snails}

The most prevalent PNW garden pest (Gardening Know How, Northwest Edible Life). The region's native banana slug---second largest slug species in the world at over 9 inches---is generally a forest-dweller. Problematic species are non-native brown and gray garden slugs (\textit{Deroceras reticulatum}, \textit{Arion circumscriptus}). Iron phosphate baits (Sluggo) are organically approved and target only slugs/snails without harming beneficial insects. Habitat modification---reducing moist debris, using raised beds---reduces populations.

\subsubsection{Imported Cabbageworm}

Larval stage of \textit{Pieris rapae}, a white butterfly with black wing spots. Particularly problematic in PNW because the climate favors both brassica crops and the pest, allowing multiple generations per summer (Northwest Edible Life). Management: floating row covers (most effective), hand-picking, encouraging parasitic wasps, and \textit{Bacillus thuringiensis} (Bt) spray.

\subsubsection{Fungal Diseases}

Powdery mildew thrives in PNW humidity and affects ornamental and edible plants widely. Prevention through disease-resistant cultivars, adequate spacing for airflow, and morning watering. Apple scab and anthracnose affect fruit and ornamental trees (Portland Bureau of Parks). Emerging threat: emerald ash borer, first detected on the West Coast in Forest Grove, Oregon in June 2022 (Portland Bureau of Parks), threatens all ash species (\textit{Fraxinus} spp.).

\subsubsection{Climate Adaptation Strategies}

The 2021 Northwest Heat Dome caused 50--100\% crop losses on many farms; the 2015 PNW drought produced an estimated \$1 billion in agricultural losses across Idaho, Oregon, and Washington (USDA Northwest Climate Hub). Hotter, drier summers are projected to continue, with longer dry seasons and more frequent extreme heat events (U.S. EPA Northwest Climate Impacts).

Garden-level adaptations: drought-resistant crop varieties; drip or soaker-hose irrigation; shade cloth for heat dome events; mesh coverings (20--30\% less crop damage in tree fruit operations per USDA Climate Hub data); no-till soil management for moisture retention; organic mulch (3 inches maximum); and strategic microclimate use (south walls for heat-loving crops, north exposures for cool-season greens in summer).

The shifting frost-free season is a double-edged opportunity: earlier planting is possible, and formerly tender perennials (certain salvias, abutilon, eucalyptus, gardenia, jasmine) may survive Zone 9a winters (Sky Nursery). However, unpredictable cold snaps remain a risk.

\subsection{Safety and Sensitivity Considerations}

\begin{center}
\rowcolors{2}{rowgray}{white}
\begin{tabularx}{\textwidth}{L{3.5cm}XC{2cm}}
\toprule
\rowcolor{primary}
\tableheader{Consideration} & \tableheader{Boundary} & \tableheader{Type} \\
\midrule
Source quality & All citations from gov agencies, university extension, or professional organizations & ABSOLUTE \\
Pesticide guidance & Only organically-approved or IPM methods; no broad-spectrum chemicals without safety context & GATE \\
Climate projections & From published USDA, NOAA, or EPA data only & GATE \\
Indigenous plant knowledge & Attributed to specific nations; no generic references & ABSOLUTE \\
Invasive species & Never recommend planting known invasive species & GATE \\
Medical claims & No health claims without peer-reviewed evidence & ABSOLUTE \\
\bottomrule
\end{tabularx}
\end{center}

\subsection{Source Bibliography}

\subsubsection{Government and Agency Sources}

\begin{itemize}[leftmargin=2em]
  \item USDA Plant Hardiness Zone Map, 2023 revision. \url{https://planthardiness.ars.usda.gov/}
  \item USDA Northwest Climate Hub --- Drought and Agriculture. \url{https://www.climatehubs.usda.gov/hubs/northwest}
  \item U.S. EPA --- Climate Impacts in the Northwest
  \item USDA NRCS --- Washington Soil Atlas (K. Hipple)
  \item Portland Bureau of Parks --- Tree Diseases and Pests
  \item King County DNRP --- Native Plant Resources for the Pacific Northwest
\end{itemize}

\subsubsection{University Extension and Research}

\begin{itemize}[leftmargin=2em]
  \item OSU Extension --- EC-1577: Gardening with Oregon Native Plants West of the Cascades (L. McMahan)
  \item OSU Extension --- Monthly Garden Calendars
  \item OSU --- Landscape Plants / USDA Hardiness Zone Maps
  \item WSU Extension --- EM051E: Home Vegetable Gardening in Washington
  \item WSU / OSU / U of Idaho --- Pacific Northwest Pest Management Handbooks
  \item WSU Mt. Vernon NWREC --- Pacific Northwest Vegetable Extension Group
  \item University of Washington --- Climate Impacts Group
\end{itemize}

\subsubsection{Professional Organizations}

\begin{itemize}[leftmargin=2em]
  \item Tilth Alliance --- Maritime Northwest Garden Guide (2014 rev. ed.)
  \item Washington Native Plant Society
  \item Pacific Horticulture Society
  \item Portland Audubon / Columbia Land Trust --- Backyard Habitat Certification
  \item Cloud Mountain Farm Center and Nursery
  \item Deep Harvest Farm --- Pacific NW Planting Calendar
\end{itemize}

\newpage

% STAGE 3 --- MISSION PACKAGE
\stageheader{STAGE 3 --- MISSION PACKAGE}{tertiary}

\section{Milestone Specification}

\subsection{Mission Objective}

Produce a comprehensive, publication-quality reference on Pacific Northwest gardening (west of the Cascades), covering climate interpretation, soil management, native plant integration, seasonal food production, and pest/disease/climate adaptation. All sources from government agencies, university extension, or professional organizations. Self-contained for a reader with no prior PNW gardening knowledge.

\subsection{Architecture Overview}

\begin{asciibox}
\begin{verbatim}
  WAVE 0          WAVE 1              WAVE 2         WAVE 3
  Foundation      Parallel Research    Synthesis       Publication
  +------+       +------+------+     +------+       +------+
  |Schema|------>|Track |Track ||--->|Cross-|------>|Final |
  |Index |       |  A   |  B   |     |Module|       |Doc   |
  |Setup |       |M1+M2 |M3+M4 |     |Synth |       |Verify|
  +------+       +------+------+     | + M5 |       +------+
                                     +------+
  Haiku           Sonnet + Opus       Opus            Sonnet
  Sequential      Parallel            Sequential      Sequential
\end{verbatim}
\end{asciibox}

\subsection{System Layers}

\begin{enumerate}[leftmargin=2em]
  \item \textbf{Foundation Layer} --- Shared schemas, source index, citation format, terminology glossary
  \item \textbf{Survey Layer} --- Five module surveys drawing from the Research Reference
  \item \textbf{Synthesis Layer} --- Cross-module integration (native plants to pest management, soil to climate adaptation)
  \item \textbf{Publication Layer} --- Final document assembly, verification, formatting
\end{enumerate}

\subsection{Deliverables}

\begin{center}
\rowcolors{2}{rowgray}{white}
\begin{tabularx}{\textwidth}{C{0.5cm}L{3.5cm}XC{2cm}}
\toprule
\rowcolor{primary}
\tableheader{\#} & \tableheader{Deliverable} & \tableheader{Acceptance Criteria} & \tableheader{Spec} \\
\midrule
1 & Source Index & All sources catalogued with quality rating & W0-SCHEMA \\
2 & Climate Reference & Zones, microclimates, seasonal patterns & M1-CLIMATE \\
3 & Soil Guide & Types, testing, amendment protocols & M2-SOIL \\
4 & Native Plant Guide & 25+ species profiles with ecology & M3-NATIVE \\
5 & Planting Calendar & 12-month, 20+ crops, region-specific & M4-FOOD \\
6 & Pest/Adaptation Guide & IPM strategies + climate resilience & M5-ADAPT \\
7 & Integrated Document & Cross-referenced, self-contained & SYNTHESIS \\
8 & Final Publication & Verified, formatted, print-ready & PUB \\
\bottomrule
\end{tabularx}
\end{center}

\subsection{Component Breakdown}

\begin{center}
\rowcolors{2}{rowgray}{white}
\begin{tabularx}{\textwidth}{L{3cm}L{2.5cm}C{1.5cm}C{2cm}}
\toprule
\rowcolor{primary}
\tableheader{Component} & \tableheader{Dependencies} & \tableheader{Model} & \tableheader{Est. Tokens} \\
\midrule
W0-SCHEMA & None & Haiku & 3,000 \\
M1-CLIMATE & W0-SCHEMA & Sonnet & 12,000 \\
M2-SOIL & W0-SCHEMA & Sonnet & 10,000 \\
M3-NATIVE & W0-SCHEMA & Opus & 15,000 \\
M4-FOOD & W0-SCHEMA, M1 & Sonnet & 14,000 \\
M5-ADAPT & W0-SCHEMA, M1 & Opus & 12,000 \\
SYNTHESIS & M1--M5 & Opus & 8,000 \\
PUB-VERIFY & SYNTHESIS & Sonnet & 5,000 \\
\midrule
\textbf{Total} & & & \textbf{79,000} \\
\bottomrule
\end{tabularx}
\end{center}

\subsection{Model Assignment Rationale}

\begin{center}
\rowcolors{2}{rowgray}{white}
\begin{tabularx}{\textwidth}{C{1.5cm}L{3cm}X}
\toprule
\rowcolor{primary}
\tableheader{Model} & \tableheader{Components} & \tableheader{Rationale} \\
\midrule
Opus & M3, M5, SYNTHESIS & Ecological integration judgment, cross-module pattern recognition, nuanced climate adaptation \\
Sonnet & M1, M2, M4, PUB & Structured survey compilation, calendar assembly, document formatting, verification \\
Haiku & W0-SCHEMA & Template scaffolding, index generation \\
\bottomrule
\end{tabularx}
\end{center}

\section{Wave Execution Plan}

\subsection{Wave Summary}

\begin{center}
\rowcolors{2}{rowgray}{white}
\begin{tabularx}{\textwidth}{C{1cm}L{3cm}C{2cm}C{1.5cm}C{2cm}}
\toprule
\rowcolor{primary}
\tableheader{Wave} & \tableheader{Name} & \tableheader{Components} & \tableheader{Mode} & \tableheader{Est. Time} \\
\midrule
0 & Foundation & W0-SCHEMA & Sequential & 1 window \\
1 & Parallel Research & M1--M5 & Parallel & 4 windows \\
2 & Synthesis & SYNTHESIS & Sequential & 1 window \\
3 & Publication & PUB-VERIFY & Sequential & 1 window \\
\bottomrule
\end{tabularx}
\end{center}

\subsection{Wave 0: Foundation}

\begin{center}
\rowcolors{2}{rowgray}{white}
\begin{tabularx}{\textwidth}{L{2.5cm}L{2cm}C{1.5cm}X}
\toprule
\rowcolor{primary}
\tableheader{Component} & \tableheader{Agent} & \tableheader{Model} & \tableheader{Output} \\
\midrule
W0-SCHEMA & PLAN & Haiku & Citation schema, source index, shared glossary, module output format \\
\bottomrule
\end{tabularx}
\end{center}

\subsection{Wave 1: Parallel Research Tracks}

\textbf{Track A: Physical Foundation (M1 + M2)}

\begin{center}
\rowcolors{2}{rowgray}{white}
\begin{tabularx}{\textwidth}{L{2.5cm}L{2.5cm}C{1.3cm}X}
\toprule
\rowcolor{primary}
\tableheader{Component} & \tableheader{Agent} & \tableheader{Model} & \tableheader{Output} \\
\midrule
M1-CLIMATE & CRAFT-CLIMATE & Sonnet & Zone guide, microclimate catalog, seasonal patterns \\
M2-SOIL & CRAFT-SOIL & Sonnet & Soil profiles, pH protocol, amendment guide, testing resources \\
\bottomrule
\end{tabularx}
\end{center}

\textbf{Track B: Biological Systems (M3 + M4)}

\begin{center}
\rowcolors{2}{rowgray}{white}
\begin{tabularx}{\textwidth}{L{2.5cm}L{2.5cm}C{1.3cm}X}
\toprule
\rowcolor{primary}
\tableheader{Component} & \tableheader{Agent} & \tableheader{Model} & \tableheader{Output} \\
\midrule
M3-NATIVE & CRAFT-FLORA & Opus & 25+ species profiles, community associations, ecology mapping \\
M4-FOOD & CRAFT-FOOD & Sonnet & 12-month calendar, 20+ crop profiles, season extension \\
\bottomrule
\end{tabularx}
\end{center}

\textbf{M5 begins after M1 completes (climate data dependency):}

\begin{center}
\rowcolors{2}{rowgray}{white}
\begin{tabularx}{\textwidth}{L{2.5cm}L{2.5cm}C{1.3cm}X}
\toprule
\rowcolor{primary}
\tableheader{Component} & \tableheader{Agent} & \tableheader{Model} & \tableheader{Output} \\
\midrule
M5-ADAPT & CRAFT-ADAPT & Opus & IPM strategies, climate adaptation, heat/drought response \\
\bottomrule
\end{tabularx}
\end{center}

\subsection{Wave 2: Synthesis}

\begin{center}
\rowcolors{2}{rowgray}{white}
\begin{tabularx}{\textwidth}{L{2.5cm}L{2cm}C{1.3cm}X}
\toprule
\rowcolor{primary}
\tableheader{Component} & \tableheader{Agent} & \tableheader{Model} & \tableheader{Output} \\
\midrule
SYNTHESIS & INTEG & Opus & Cross-module integration: natives as pest habitat, soil as climate buffer, calendar linked to zones \\
\bottomrule
\end{tabularx}
\end{center}

\subsection{Wave 3: Publication and Verification}

\begin{center}
\rowcolors{2}{rowgray}{white}
\begin{tabularx}{\textwidth}{L{2.5cm}L{2cm}C{1.3cm}X}
\toprule
\rowcolor{primary}
\tableheader{Component} & \tableheader{Agent} & \tableheader{Model} & \tableheader{Output} \\
\midrule
PUB-VERIFY & VERIFY & Sonnet & Source validation, accuracy check, self-containment test, formatting \\
\bottomrule
\end{tabularx}
\end{center}

\subsection{Cache Optimization Strategy}

\begin{itemize}[leftmargin=2em]
  \item Wave 0 output (schemas, templates) is lightweight; load into all Wave 1 agent contexts
  \item Track A and Track B run in parallel within the 5-minute cache TTL window
  \item M5 begins as M1 completes, reusing cached climate context
  \item Synthesis receives all five module outputs; module summaries (first 500 tokens each) can substitute if context limits require
  \item Final verification needs only the synthesized document plus Wave 0 source index
\end{itemize}

\subsection{Token Budget}

\begin{center}
\rowcolors{2}{rowgray}{white}
\begin{tabularx}{\textwidth}{L{3cm}C{2cm}C{2cm}C{2cm}}
\toprule
\rowcolor{primary}
\tableheader{Model} & \tableheader{Components} & \tableheader{Est. Tokens} & \tableheader{Budget \%} \\
\midrule
Opus & M3, M5, SYNTH & 35,000 & 44\% \\
Sonnet & M1, M2, M4, PUB & 41,000 & 52\% \\
Haiku & W0-SCHEMA & 3,000 & 4\% \\
\midrule
\textbf{Total} & & \textbf{79,000} & \textbf{100\%} \\
\bottomrule
\end{tabularx}
\end{center}

\subsection{Dependency Graph}

\begin{asciibox}
\begin{verbatim}
  W0-SCHEMA
      |
      +----------------------------------+
      |                                  |
      v                                  v
  +--------+  +--------+          +--------+  +--------+
  |M1-CLIM |  |M2-SOIL |          |M3-NATV |  |M4-FOOD |
  +----+---+  +--------+          +--------+  +--------+
       |
       v
  +--------+
  |M5-ADPT |
  +--------+
       |
       +---- all M1-M5 complete ----+
       |                            |
       v                            v
  +---------+                 +----------+
  |SYNTHESIS|---------------->|PUB-VERIFY|
  +---------+                 +----------+
\end{verbatim}
\end{asciibox}

\section{Test Plan}

\subsection{Test Categories}

\begin{center}
\rowcolors{2}{rowgray}{white}
\begin{tabularx}{\textwidth}{L{3cm}C{1.5cm}C{1.5cm}C{2cm}C{2cm}}
\toprule
\rowcolor{primary}
\tableheader{Category} & \tableheader{Count} & \tableheader{Priority} & \tableheader{Failure} & \tableheader{Target \%} \\
\midrule
Safety-critical & 7 & Mandatory & BLOCK & 18\% \\
Core functionality & 17 & Required & BLOCK & 45\% \\
Integration & 9 & Required & BLOCK & 24\% \\
Edge cases & 5 & Best-effort & LOG & 13\% \\
\midrule
\textbf{Total} & \textbf{38} & & & \textbf{100\%} \\
\bottomrule
\end{tabularx}
\end{center}

\subsection{Safety-Critical Tests}

\begin{center}
\rowcolors{2}{rowgray}{white}
\begin{tabularx}{\textwidth}{C{1.2cm}L{3.5cm}X}
\toprule
\rowcolor{primary}
\tableheader{ID} & \tableheader{Verifies} & \tableheader{Expected Behavior} \\
\midrule
SC-SRC & Source quality & All citations from gov agencies, extension, or professional orgs. Zero entertainment media. \\
SC-NUM & Numerical attribution & Every temperature, pH, zone, species count, and economic figure attributed to a specific source. \\
SC-ADV & No policy advocacy & Evidence presented without advocating specific legislative positions. \\
SC-CLI & Climate projections & Every projection cites USDA, NOAA, EPA, or IPCC data. \\
SC-INV & Invasive species & No recommendations include species listed as invasive in OR or WA. \\
SC-MED & Medical/health claims & No unattributed health claims for edible or medicinal plants. \\
SC-PST & Pesticide safety & Every pesticide recommendation includes organic status and non-target impact. \\
\bottomrule
\end{tabularx}
\end{center}

\subsection{Core Functionality Tests (Selected)}

\begin{center}
\rowcolors{2}{rowgray}{white}
\begin{tabularx}{\textwidth}{C{1.2cm}L{3.5cm}X}
\toprule
\rowcolor{primary}
\tableheader{ID} & \tableheader{Verifies} & \tableheader{Expected Behavior} \\
\midrule
CF-01 & Zone coverage & Climate module covers USDA zones 7b--9a with practical interpretation \\
CF-02 & Sunset zones & At least Sunset zones 4--6 described \\
CF-03 & Soil type coverage & All four major types profiled with amendments \\
CF-04 & pH management & Lime thresholds and testing resources cited \\
CF-05 & Native plant count & Minimum 25 species across all categories \\
CF-06 & Species profiles & Each species has zone, sun/moisture, ecological function \\
CF-07 & Calendar breadth & 12-month, at least 20 vegetables with timing \\
CF-08 & Succession planting & At least 3 crops with succession schedule \\
CF-09 & Top 5 pests & Slugs, cabbageworm, weevils, earwigs, and one disease covered \\
CF-10 & IPM strategies & Each pest has at least 2 organic/ecological management options \\
CF-11 & Drought adaptation & At least 4 garden-level strategies documented \\
CF-12 & Heat dome response & Shade cloth, cooling, microclimate strategies documented \\
\bottomrule
\end{tabularx}
\end{center}

\subsection{Integration Tests (Selected)}

\begin{center}
\rowcolors{2}{rowgray}{white}
\begin{tabularx}{\textwidth}{C{1.2cm}L{3.5cm}X}
\toprule
\rowcolor{primary}
\tableheader{ID} & \tableheader{Verifies} & \tableheader{Expected Behavior} \\
\midrule
IN-01 & Native to Pest cascade & Traces native habitat to beneficial insects to pest suppression \\
IN-02 & Soil to Climate cascade & No-till and mulching linked to soil health and climate adaptation \\
IN-03 & Calendar to Climate & Planting dates consistent with frost dates and zone data \\
IN-04 & Cross-module consistency & Same species referenced consistently across modules \\
IN-05 & Self-containment & No undefined terms, complete bibliography \\
\bottomrule
\end{tabularx}
\end{center}

\subsection{Verification Matrix}

\begin{center}
\rowcolors{2}{rowgray}{white}
\begin{tabularx}{\textwidth}{XL{3cm}C{1.5cm}}
\toprule
\rowcolor{primary}
\tableheader{Success Criterion} & \tableheader{Test ID(s)} & \tableheader{Status} \\
\midrule
1. Zone systems covered & CF-01, CF-02 & Pending \\
2. Soil types with amendments & CF-03, CF-04 & Pending \\
3. 25+ native species profiled & CF-05, CF-06 & Pending \\
4. 12-month, 20+ crop calendar & CF-07, CF-08 & Pending \\
5. Top 5 pests with IPM & CF-09, CF-10 & Pending \\
6. Climate projections sourced & SC-CLI & Pending \\
7. Plant recs include zone/sun/moisture & CF-06 & Pending \\
8. Climate adaptation strategies & CF-11, CF-12 & Pending \\
9. Cross-module integration & IN-01, IN-02 & Pending \\
10. All numbers attributed & SC-NUM & Pending \\
11. Self-contained document & IN-05 & Pending \\
12. No policy advocacy & SC-ADV & Pending \\
\bottomrule
\end{tabularx}
\end{center}

\section{Mission Crew Manifest}

\begin{center}
\rowcolors{2}{rowgray}{white}
\begin{tabularx}{\textwidth}{L{2cm}L{3cm}X}
\toprule
\rowcolor{primary}
\tableheader{Role} & \tableheader{Agent} & \tableheader{Responsibility} \\
\midrule
FLIGHT & Orchestrator & Mission direction; wave go/no-go gates \\
PLAN & Planner & Wave decomposition; parallel track optimization; cache strategy \\
SCOUT & Research Agent & Source gathering; web research for gap-fill \\
EXEC-A & Executor (Track A) & M1 and M2 document production \\
EXEC-B & Executor (Track B) & M3 and M4 document production \\
CRAFT-CLIM & Climate Specialist & Zone interpretation, microclimate analysis \\
CRAFT-FLORA & Botany Specialist & Native plant profiles, community ecology \\
VERIFY & Verification Agent & Source validation; numerical accuracy; completeness \\
INTEG & Integration Agent & Cross-module relationship mapping and synthesis \\
CAPCOM & HITL Interface & Human approval at wave boundaries \\
BUDGET & Resource Manager & Token budget tracking; session planning \\
LOG & Chronicle Agent & Commit messages; milestone audit trail \\
\bottomrule
\end{tabularx}
\end{center}

\section{Execution Summary}

\begin{center}
\rowcolors{2}{rowgray}{white}
\begin{tabularx}{\textwidth}{L{5cm}X}
\toprule
\rowcolor{primary}
\tableheader{Metric} & \tableheader{Value} \\
\midrule
Activation Profile & Squadron (12 roles) \\
Total Modules & 5 research + 1 synthesis + 1 publication \\
Parallel Tracks & 2 (Physical Foundation + Biological Systems) \\
Total Waves & 4 (Foundation, Research, Synthesis, Publication) \\
Estimated Context Windows & 12 \\
Estimated Sessions & 4 \\
Total Token Budget & 79,000 \\
Model Split & Opus 44\% / Sonnet 52\% / Haiku 4\% \\
Test Count & 38 (7 safety-critical, 17 core, 9 integration, 5 edge) \\
Safety-Critical Tests & 7 (all mandatory-pass BLOCK) \\
\bottomrule
\end{tabularx}
\end{center}

\vspace{2em}

\begin{quotebox}
\textit{``The garden is not a place apart from the world. It is the world at the scale where we can see the connections---between rain and root, between flower and wing, between the patience of soil and the ambition of seed. To garden in the Pacific Northwest is to join a conversation that has been going on for ten thousand years.''} \par\vspace{0.5em}
\hfill --- \textbf{The Through-Line: where architecture meets ecology}
\end{quotebox}

\end{document}
